%(BEGIN_QUESTION)
% Copyright 2006, Tony R. Kuphaldt, released under the Creative Commons Attribution License (v 1.0)
% This means you may do almost anything with this work of mine, so long as you give me proper credit

John har en sykkel med hjuldiameter på 26 tommer. Hvor fort roterer hjulene (i omdreininger per minutt, «RPM») når John sykler i 15 miles per hour? Husk at omkretsen av en sirkel er lik diameteren multiplisert med $\pi$. Denne verdien er strekningen hjulet tilbakelegger per omdreining.

\underbar{file i00105}
%(END_QUESTION)





%(BEGIN_ANSWER)

15 MPH = 193.924 RPM

\vskip 10pt

Ved første øyekast virker dette kanskje ikke som en enhetsomregning, men det er det. Hjulene står selvsagt stille ved 0 miles per hour, og rotasjonshastigheten følger sykkelhastigheten lineært. Det vi har med å gjøre er en direkte proporsjon, og det er nettopp dette enhetsomregninger handler om.

Nøkkelen til oppgaven er å lage en enhetsbrøk som knytter én hjulomdreining til tilbakeelagt strekning. Holder vi oss i tommer, får vi 81.68 tommer per omdreining. Disse tallene bruker vi for å bygge enhetsbrøken som konverterer hastighet til rotasjonshastighet (fjerde brøk i denne lange utregningen):

$${15 \hbox{ mi} \over 1 \hbox{ t}} \times {5280 \hbox{ ft} \over 1 \hbox{ mi}} \times {12 \hbox{ in} \over 1 \hbox{ ft}} \times {1 \hbox{ rev} \over 81.68 \hbox{ in}} \times {1 \hbox{ t} \over 60 \hbox{ min}}$$

%(END_ANSWER)





%(BEGIN_NOTES)


%INDEX% Physics, units and conversions: unit conversions using unity fractions

%(END_NOTES)
